%% ======================================================================
\section{Posteriors, priors and the computation of the bounds}
\label{chap:posteriors}
%% ======================================================================

\begin{objectives}
\guiding{Since the bounds hold for every prior and every posterior, which ones
should we choose, and how do we evaluate the bound in practice?}
The bounds of \cref{chap:bounds} hold for any prior and any posterior. This
section explains how to choose them. We show that the posterior minimizing a
linear bound is the Gibbs posterior, a soft selection of the good voters that
also contains Bayesian inference as a special case, and we run it on a real
dataset. We then discuss the choice of the prior, which is the main lever to
obtain tight certificates, and we measure on an experiment how to split the data
between the prior and the bound. We finally explain how to compute a bound
when the Gibbs risk has no closed form, and how PAC-Bayes can say something about
deterministic predictors.
\end{objectives}

\subsection{The optimal posterior: Gibbs posterior and Bayesian inference}
\label{sec:gibbs_posterior}

A bound that holds for every posterior invites an obvious question: which
posterior makes it smallest? For linear bounds the answer is explicit. The
bounds \eqref{eq:catoni_hoeffding} and \eqref{eq:lambda} depend on
$Q$ only through the objective
\begin{equation}\label{eq:objective_linear}
\RS(\GQ)+\frac{\KL(Q\|P)}{\lambda m}
=\Esp{h\sim Q}{\RS(h)}+\frac{1}{\lambda m}\KL(Q\|P),
\end{equation}
which balances the fit of the posterior to the data against its distance to the
prior; the parameter $\lambda$ plays the role of an inverse temperature. We
already know how to minimize such an objective: this is exactly the variational formula
\eqref{eq:gibbs_variational} of \cref{chap:tools}.

\begin{proposition}{Gibbs posterior \citep{catoni2007pac,zhang2006information}}{gibbs_post}
The minimizer over $\M(\Hcal)$ of \eqref{eq:objective_linear} is the
\textbf{Gibbs posterior}
\begin{equation}\label{eq:gibbs_post}
Q_\lambda(h)=\frac{P(h)\,e^{-\lambda m\RS(h)}}{\Esp{h'\sim P}{e^{-\lambda m\RS(h')}}},
\end{equation}
and the minimal value is $-\frac{1}{\lambda m}\ln\E_{h\sim P}\,e^{-\lambda m\RS(h)}$.
\end{proposition}
\begin{proof}
Apply the variational formula \eqref{eq:gibbs_variational} to
$\phi(h)=-\lambda m\RS(h)$: $\sup_Q\{-\lambda m\E_Q\RS-\KL(Q\|P)\}=\ln\E_Pe^{-\lambda m\RS}$,
with the supremum attained at $Q_\phi=Q_\lambda$. Multiply by $-\frac1{\lambda m}$.
\end{proof}

In words, the Gibbs posterior starts from the prior and reweights it
exponentially according to the empirical risk: a voter that makes $k$ more
training errors than another sees its relative weight divided by $e^{\lambda k}$.
The temperature controls how sharp this selection is. When $\lambda\to0$, the
posterior stays equal to the prior: there is no complexity to pay, but nothing
is learned. When
$\lambda\to\infty$, the posterior concentrates on the empirical risk minimizers;
for a finite class with the uniform prior, its KL tends to $\ln|\Hcal|$ minus
the logarithm of the number of minimizers, and we recover the Occam bound.
\Cref{fig:prior_posterior} shows this transition on a family of threshold
classifiers $h_\theta(x)=\sign(x-\theta)$ with a Gaussian prior on $\theta$: as
$\lambda$ grows, the posterior moves from the prior towards the interval of
thresholds that minimize the empirical risk.

\begin{figure}[t]
\centering
\includegraphics[width=\linewidth]{prior_posterior}
\caption{Gibbs posteriors on a family of threshold classifiers
$h_\theta(x)=\sign(x-\theta)$ with a Gaussian prior on $\theta$. As
$\lambda$ increases, the posterior concentrates on the empirical risk minimizers;
for small $\lambda$ it stays close to the prior.}
\label{fig:prior_posterior}
\end{figure}

\begin{example}[The Gibbs posterior on decision stumps]
Let us see the temperature at work on the \texttt{breast cancer} dataset. We build
$300$ decision stumps (one feature, one threshold, one orientation). Their
thresholds are the quantiles of the features computed on $100$ examples that
are then set aside, so that the set of voters, and hence the uniform prior, is
independent of the $m=120$ examples on which the bound is computed. We then
compute $Q_\lambda$ for $80$ values of $\lambda$ between $10^{-3}$ and $10^2$,
and Seeger's bound for each of them.

The right panel of \cref{fig:temperature} shows the trade-off. For small
$\lambda$, the posterior is uniform, its empirical Gibbs risk is $0.5$ (each
stump comes with its opposite), and the bound is vacuous. As $\lambda$ grows,
the empirical risk drops quickly while the KL increases slowly. The bound
reaches its minimum $0.256$ for $\lambda\approx1.3$, where the empirical Gibbs
risk is $0.084$ and $\KL(Q_\lambda\|P)=5.5$ nats; the Gibbs risk measured on the
$349$ test examples is $0.114$, well below the certificate. Beyond this point,
the KL keeps increasing towards $\ln 300\approx5.7$ while the empirical risk no
longer decreases.
\end{example}

\begin{figure}[t]
\centering
\includegraphics[width=\linewidth]{gibbs_temperature}
\caption{Gibbs posterior over $300$ decision stumps on the \texttt{breast cancer}
dataset (thresholds at quantiles computed on $100$ separate examples, uniform
prior, bound computed on $m=120$ other examples). Left: posterior weights
for three temperatures. Right: as $\lambda$ grows, the empirical Gibbs risk
decreases while $\KL(Q_\lambda\|P)$ increases, up to $\ln|\Hcal|$ (Occam);
Seeger's bound is minimized for an intermediate temperature.}
\label{fig:temperature}
\end{figure}

With only $m=120$ examples, a certificate of $0.256$ may seem loose. It is
however a \emph{guarantee}, valid with probability $0.95$ and obtained without
any test set. The example also shows that the Gibbs posterior is essentially as
good as the best stump in terms of certificate, while being more stable. What
it certifies, though, is a stump drawn at random; we will see in \cref{chap:mv}
how to certify the \emph{vote} of these stumps instead.

\begin{running}
The same recipe applies to the $100$ trees of our forest, with their
out-of-bag errors in place of the empirical risks. Choosing the temperature by
the procedure of \cref{sec:algo_fo}, the Gibbs posterior lowers Seeger's
certificate on the Gibbs risk from $0.230$ (uniform weights) to $0.193$. This
looks like progress, but there is a catch, which we will only understand in
\cref{chap:mv}: the Gibbs posterior puts almost all its weight on the two best
trees, and the test risk of the \emph{vote} jumps from $0.044$ to $0.150$. A
better Gibbs classifier does not always make a better vote.
\end{running}

\paragraph{Link with Bayesian inference.}
The Gibbs posterior is not only an algorithmic device: it contains Bayesian
inference as a special case. Take a probabilistic model $p(y\,|\,x,\theta)$
indexed by $\theta\in\Hcal$, and the (unbounded) negative log-likelihood loss
$\ell(\theta,(x,y))=-\ln p(y\,|\,x,\theta)$. For $\lambda=1$, the Gibbs posterior is
\[
Q_1(\theta)\propto P(\theta)\,e^{-\sum_{i=1}^m\ell(\theta,(x_i,y_i))}=P(\theta)\prod_{i=1}^mp(y_i\,|\,x_i,\theta),
\]
which is exactly the Bayesian posterior. \citet{germain2016pac} exploited this
correspondence to obtain PAC-Bayesian guarantees for Bayesian inference, and to
show that minimizing the bound amounts to maximizing the \emph{marginal
likelihood} $\E_{\theta\sim P}\prod_i p(y_i|x_i,\theta)$, which justifies model
selection by the evidence. Since the log-loss is unbounded, the bound requires
additional assumptions on the tails of the loss (sub-Gaussian or sub-gamma).
Other values of $\lambda$ give \emph{tempered}, or ``generalized'', posteriors,
and \citet{alquier2016properties} study variational approximations of Gibbs
posteriors when they cannot be computed exactly.

The formulas coincide, but the two frameworks read them differently. In the
Bayesian view the prior encodes beliefs
and the posterior is the result of conditioning, whereas in PAC-Bayes the prior
is a reference measure, the posterior is any distribution we like, and the
guarantee does not require the model to be correct.

\begin{intuition}
PAC-Bayes explains \emph{why} averaging with exponential weights is a good idea:
it is the optimal way to trade empirical risk against a KL penalty. In a
finite ensemble (\cref{chap:algos}), the Gibbs posterior simply assigns to
each voter a weight proportional to $P(h)\,e^{-\lambda m\RS(h)}$, a soft version
of ``keep the best voters''.
\end{intuition}

\subsection{Choosing the prior}\label{sec:prior}

The Gibbs posterior tells us how to choose $Q$ once $P$ is given. But the
prior matters just as much, since the whole complexity term is measured
relative to it. The only rule is that the prior must not depend on the sample
used to evaluate the bound; otherwise it can be anything. Since the bound pays
$\KL(Q\|P)$, it is tight only if the prior already puts mass where the
posterior will be. Much of the art of PAC-Bayes thus consists in finding priors
that are independent of the sample and yet well informed. We review the main
strategies, from the most naive to the most data-efficient.

The simplest priors are \emph{data-independent}: the uniform distribution on a
finite set of voters, a Gaussian $\mathcal N(0,\sigma^2 I)$ on the weights of
linear classifiers or neural networks, or the Fourier transform of a kernel for
random features \citep{letarte2019pseudo}. Their parameters, such as the scale
$\sigma$, can be selected on a grid at the price of a union bound. They are
often too vague, however: nothing tells them in which region of $\Hcal$ the
good hypotheses lie.

To make the prior better informed without breaking the rules, a more effective strategy is to learn the prior on a part of the data that is
not used in the bound \citep{ambroladze2006tighter,parrado2012pac}. The sample is
split into $S=S_{\text{prior}}\cup S_{\text{bound}}$; the prior is learned on
$S_{\text{prior}}$, for instance a Gaussian centred at an SVM trained on it,
and the bound is computed on $S_{\text{bound}}$ only, with $m=|S_{\text{bound}}|$.
This is the approach behind the tightest certificates known for neural networks
\citep{dziugaite2017computing,perez2021tighter}. The split involves a trade-off:
a larger $S_{\text{prior}}$ gives a better prior, hence a smaller KL, but it
leaves fewer examples to compute the bound. How should we balance the two?

\begin{example}[How much data for the prior?]
We measure this trade-off on \texttt{breast cancer} (\cref{fig:ddprior}). A
fraction of the data is used to build a set of $200$ shallow decision trees on
bootstrap samples. The uniform distribution over these trees plays the role of
the prior, which is legitimate since it only depends on $S_{\text{prior}}$. The
remaining examples are split into a sample $S_{\text{bound}}$ ($70\%$), on which
we minimize Seeger's bound over Gibbs posteriors, and a test sample ($30\%$).

When only $5\%$ of the data are used for the prior, the trees are weak, but the
bound is computed on $378$ examples and the certificate is $0.142$. Increasing
the fraction to $10$--$30\%$ improves the trees (the test Gibbs risk decreases
from $0.080$ to $0.067$) and the certificate stays around $0.135$. Beyond $50\%$,
the trees keep improving but the bound is computed on too few examples: with
$85\%$ of the data for the prior, only $60$ examples remain, and the certificate
degrades to $0.225$ although the voters are the best ones. On this dataset, a
fifth to a third of the data for the prior is a good compromise. The best
proportion depends on the problem; like any hyperparameter of the certificate,
it must be fixed in advance or selected on a grid with a union bound.
\end{example}

\begin{figure}[t]
\centering
\includegraphics[width=0.72\linewidth]{ddprior_split}
\caption{Splitting the data between the construction of the voters (prior) and
the computation of the bound, on \texttt{breast cancer} (mean and standard
deviation over $5$ splits; $m$ is the size of the sample used by the bound).}
\label{fig:ddprior}
\end{figure}

Any split, however well chosen, sacrifices part of the sample. Two further ideas
avoid this waste altogether. In bagging, each voter is
trained on a subsample and can be evaluated on its out-of-bag examples; this
\emph{split construction} lets every example serve both for building some
voters and for evaluating the others \citep{thiemann2017strongly,lorenzen2019pac},
and it is the backbone of the self-bounding algorithms of \cref{sec:oob}.
Finally, the prior may depend on the distribution $\D$ without depending on the
sample, for instance $P\propto e^{-\beta\RD(h)}$; the bound then involves unknown
quantities which can be controlled \citep{catoni2007pac}. \citet{dziugaite2018data}
go one step further and learn the prior from $S$ itself, using differential
privacy to limit how much it depends on each example.

Once the prior and the posterior are chosen, one practical step remains:
evaluating the bound numerically.

\subsection{Computing the bound in practice}

A certificate is only useful if we can compute it. When $\Hcal$ is a finite ensemble of $n$ voters, $\RS(\GQ)=\sum_hQ(h)\RS(h)$ and
$\KL(Q\|P)$ are finite sums, and computing a bound only requires one kl
inversion. When $\Hcal$ is infinite, for instance with a Gaussian posterior on the
weights of a neural network, the empirical Gibbs risk is an integral that may
not have a closed form. \citet{langford2001not} and \citet{dziugaite2017computing}
estimate it by Monte Carlo: they draw $h_1,\dots,h_N\sim Q$ independently and
apply a second kl inversion. This is legitimate because, conditionally on $S$,
the $\RS(h_j)$ are i.i.d.\ variables in $[0,1]$ with mean $\RS(\GQ)$. With
probability at least $1-\delta'$ over the draws,
\[
\RS(\GQ)\leq\klup\Big(\frac1N\sum_{j=1}^N\RS(h_j),\ \frac{\ln(2/\delta')}{N}\Big),
\]
and plugging this upper bound into Seeger's bound gives a certificate valid with
probability at least $1-\delta-\delta'$. In other words, sampling costs us a
little confidence and a slightly larger bound, but never the validity of the
certificate.

\begin{example}[Certifying a Gaussian posterior by sampling]
How large is this extra cost? We take a linear classifier $\mathbf w$ learned by a linear SVM on
\texttt{breast cancer} ($m=398$ training examples, standardized features),
rescaled to $\|\mathbf w\|=8$, and the posterior $Q=\mathcal N(\mathbf w,I)$
with prior $P=\mathcal N(\mathbf 0,I)$, so that $\KL(Q\|P)=\|\mathbf w\|^2/2=32$. For this
family, the empirical Gibbs risk has a closed form (\cref{sec:linear}), equal to
$0.076$, and Seeger's bound gives $0.245$. \Cref{fig:montecarlo} shows the
certificate obtained when the Gibbs risk is instead estimated from $N$
sampled classifiers, with $\delta=\delta'=0.025$. With $N=100$ samples, the
certificate is $0.375$, and it reaches $0.26$ for $N=10^4$: the price of the
Monte Carlo estimation decreases quickly with $N$, and a few thousand samples are
enough in practice.
\end{example}

\begin{figure}[t]
\centering
\includegraphics[width=0.66\linewidth]{montecarlo}
\caption{Certificate of a Gaussian posterior on linear classifiers
(\texttt{breast cancer}), when the empirical Gibbs risk is computed exactly
(dashed) or estimated with $N$ sampled classifiers and a second kl inversion.}
\label{fig:montecarlo}
\end{figure}

\subsection{From stochastic to deterministic predictors}

Everything so far certifies a \emph{stochastic} predictor. A PAC-Bayesian
bound controls the \emph{average} risk of a randomly drawn voter, whereas in
practice we usually deploy a single deterministic predictor. There are three
ways to bridge this gap. The first one, which is the object of \cref{chap:mv},
is to bound the risk of the majority vote $\BQ$ using the Gibbs risk and
second-order quantities describing the correlations between voters. The second
one exploits a margin: if $Q$ is concentrated around a hypothesis $h$ with
a large margin, then $h$ and the Gibbs classifier make similar predictions, and
the bound on the latter transfers to the former up to a margin term
\citep{langford2002pac,mcallester2003simplified,biggs2022margins}. The third one
consists in \emph{disintegrated} bounds, which hold with high probability for a
single hypothesis $h\sim Q$ drawn once, rather than on average over $Q$
\citep{blanchard2007occam,rivasplata2020pac,viallard2024general}; we come back to
them in \cref{chap:advances}. For ensemble methods the first route is the
natural one, and it is the subject of the next section.

\begin{keypoints}
Linear bounds are minimized by the Gibbs posterior $Q_\lambda\propto P\,e^{-\lambda m\RS}$,
a soft selection of the good voters whose temperature trades fit against
complexity, and which contains the Bayesian posterior as a special case. The
prior is the main lever for tight bounds: it can be learned from data, provided
those data are not used by the bound, and on our example a fifth to a third of the
sample was a good investment. When the Gibbs risk is not computable exactly, a
Monte Carlo estimate and a second kl inversion give a valid certificate at a
small extra cost.
\end{keypoints}

\begin{quickchecks}
\item What happens to the Gibbs posterior when $\lambda\to0$, and when $\lambda\to\infty$?
\item Why can a prior learned on part of the data be legitimate, and on which part must the bound then be computed?
\item Why does a Monte Carlo estimate of the Gibbs risk require a second kl inversion?
\item Name three ways to obtain a guarantee for a deterministic predictor from a PAC-Bayesian bound.
\end{quickchecks}

\begin{exercises}
\begin{exercise}
For a finite class and a uniform prior, show that $\KL(Q_\lambda\|P)\to\ln\frac{|\Hcal|}{k}$ when
$\lambda\to\infty$, where $k$ is the number of empirical risk minimizers.
\end{exercise}
\begin{exercise}
Explain why using the same sample to learn the centre of a Gaussian prior and to
compute the empirical term of the bound is not allowed. What happens if we do it anyway?
\end{exercise}
\begin{exercise}
In the Monte Carlo certificate, how should $\delta$ and $\delta'$ be split to
minimize the final bound when $N=100$? when $N=10^4$?
\end{exercise}
\begin{exercise}
\textbf{(Comprehension)} (a)~What is the Gibbs posterior $Q_\lambda$ when all the voters
have the same empirical risk? (b)~With the uniform prior, voter $h$ makes $3$ more
training errors than voter $h'$; what is the ratio $Q_\lambda(h)/Q_\lambda(h')$? (c)~True or
false: a prior $P\propto e^{-\beta\RD(h)}$, which depends on $\D$, is forbidden.
\end{exercise}
\begin{exercise}
\textbf{(Application)} Four voters have empirical risks $0.10$, $0.12$, $0.15$ and $0.30$ on
$m=100$ examples, and the prior is uniform. For $\lambda\in\{0.1,0.5,2\}$, compute the Gibbs
posterior $Q_\lambda$, its empirical Gibbs risk, $\KL(Q_\lambda\|P)$ and Seeger's bound with
$\delta=0.05$. For $\lambda=0.5$, check numerically that the objective
\eqref{eq:objective_linear} at $Q_\lambda$ equals $-\frac1{\lambda m}\ln\E_{h\sim P}e^{-\lambda m\RS(h)}$.
Compare with the bounds of the uniform posterior and of the Dirac on the best voter.
\end{exercise}
\begin{exercise}
\textbf{(Application)} A practitioner has $500$ examples and $100$ voters. Option~A uses the
$500$ examples for the bound with a data-independent uniform prior; the best posterior found
has an empirical Gibbs risk of $0.15$ and $\KL(Q\|P)=2$. Option~B learns a prior on $150$
examples and computes the bound on the remaining $350$; the best posterior has an
empirical Gibbs risk of $0.10$ and $\KL(Q\|P)=1$. Compare Seeger's bounds ($\delta=0.05$).
Can the practitioner compute both and report the smaller one?
\end{exercise}
\begin{exercise}
\textbf{(Application)} In the setting of the Gaussian posterior example ($m=398$,
$\KL(Q\|P)=32$), suppose that $N=1000$ classifiers sampled from $Q$ have an average
empirical risk of $0.08$. With $\delta=\delta'=0.025$, compute the Monte Carlo upper bound on
$\RS(\GQ)$ and the final certificate. Compare with the certificate that would be obtained
if $\RS(\GQ)=0.08$ were known exactly.
\end{exercise}
\begin{exercise}
\textbf{(Proof)} Prove the formula
$\KL\big(\mathcal N(\mu_1,\sigma^2 I_d)\,\|\,\mathcal N(\mu_0,\sigma^2 I_d)\big)=\frac{\|\mu_1-\mu_0\|^2}{2\sigma^2}$,
and deduce $\KL(Q_{\mathbf w}\|P)=\frac12\|\mathbf w\|^2$ for $Q_{\mathbf w}=\mathcal N(\mathbf w,I_d)$ and
$P=\mathcal N(\mathbf 0,I_d)$.
\end{exercise}
\begin{exercise}
\textbf{(Proof)} Let $P$ be a prior on a finite set of parameters $\theta$ and
$\ell(\theta,(x,y))=-\ln p(y\,|\,x,\theta)$. Using Proposition~\ref{prop:gibbs_post} with $\lambda=1$,
show that the Bayesian posterior minimizes $\E_{\theta\sim Q}\sum_i\ell(\theta,(x_i,y_i))+\KL(Q\|P)$
over $Q$, and that the minimal value is minus the log marginal likelihood
$-\ln\E_{\theta\sim P}\prod_ip(y_i\,|\,x_i,\theta)$.
\end{exercise}
\end{exercises}
